\section{Key Results}

\paragraph{Human evaluation}
\subparagraph{Expert blind assessment}

The expert blind test provides a differentiated assessment of the two system versions (V1 and V2). While the questionnaire-based ratings reveal statistically significant differences across specific use cases and evaluation dimensions, the accompanying think-aloud data offer additional insight into how experts perceived the qualitative differences between both systems. One expert interview was excluded from the analysis because essential course vectors were isolated as technical outliers in the vector database and were therefore not retrievable, which would have distorted comparability with the remaining cases.

With regard to the use-case-level evaluation, the results show a mixed but overall favorable pattern for V2. In the technical support use case, V2 performed significantly worse than V1 ($p = 0.002$), indicating that the earlier system version was judged more positively for administrative and fact-oriented requests. In contrast, V2 achieved a slight but significant improvement in course-related subject questions ($p = 0.036$). The strongest advantage for V2 emerged in the learning assistance use case, where it was rated significantly better than V1 ($p < 0.001$). This suggests that the revised architecture was particularly beneficial for prompts requiring didactic guidance and learning-oriented interaction rather than simple factual retrieval.

A more fine-grained analysis by evaluation dimension further clarifies these differences. V2 achieved significantly higher ratings in learning support/didactics ($p < 0.01$) and relevance ($p < 0.01$). By contrast, content correctness and clarity/structure did not differ significantly between the two versions, indicating a broadly comparable level of perceived factual quality and comprehensibility. For completeness/efficiency, V1 showed a slight but non-significant advantage. Context integration/coherence showed a mean advantage for V2, although this difference did not reach statistical significance and displayed greater variance across cases. Overall, the cleaned aggregate comparison still indicates a slight but statistically significant advantage for V2.

The qualitative think-aloud analysis helps explain these quantitative findings. Experts consistently described V2 as more supportive of learning and in some cases explicitly referred to it as a ``Socratic learning assistant.'' In particular, V2 was perceived as better at prompting reflection and incorporating metacognitive elements, which is highly relevant for educational chatbot design. Experts also evaluated V2 more positively with regard to UX and transparency, especially due to improved link presentation and visible streaming behavior. The streaming output was interpreted as a useful signal that the system was still processing the request, which improved the perceived interaction quality.

At the same time, the think-aloud comments also revealed important limitations of V2. In terms of answer quality, experts characterized V1 as more likely to provide concise ``high-level'' answers, whereas V2 was perceived as integrating course-specific context more deeply. This difference was described primarily as one of abstraction level and contextual depth rather than simple quality superiority. Moreover, V2 received some critical comments regarding language consistency, especially when English prompt structures appeared in otherwise German interactions, which led to perceived inconsistency and uncertainty. This indicates that linguistic coherence is a sensitive quality criterion in educational settings that may not be fully captured by structured rating scales alone. Regarding references, experts considered both systems largely comparable, with no clear advantage for either version. Finally, latency was consistently perceived in favor of V1, which was judged to respond faster.

Taken together, the expert blind test shows that V2 did not improve all aspects of performance equally, but it did yield meaningful gains in those dimensions most closely related to the paper’s educational objective: relevance, didactic quality, and learning-oriented interaction. The results therefore suggest that the best-practice-inspired redesign particularly strengthens the chatbot’s role as a learning assistant, while trade-offs remain in technical support performance, language consistency, and perceived response speed.

